%==========================================================
% CIKM 2026 — Abstract draft (internal review)
% Trajectory-Anchored Dynamic Data Selection for Instruction Tuning
%==========================================================
\documentclass[sigconf,nonacm]{acmart}

\let\Bbbk\relax
\usepackage{amssymb}
\usepackage{amsthm}

%-------- ACM conference metadata --------
\acmConference[CIKM '26]{The 35th ACM International Conference on
                        Information and Knowledge Management}
                       {November 7--11, 2026}{Rome, Italy}
\acmYear{2026}
\copyrightyear{2026}

\settopmatter{printacmref=false, printccs=false, printfolios=true}
\pagestyle{plain}

\renewcommand{\shortauthors}{Shin et al.}

\begin{document}

\title{Trajectory-Anchored Dynamic Data Selection
       for Instruction Tuning}

%==========================================================
% Authors (one-line names, shared affiliation, 4 emails)
%==========================================================
\author{Jieun Shin, Minsoo Kim, Kyunglim Kim, and James Russell Geraci}
\authornote{James Russell Geraci is the corresponding author.}
\affiliation{%
  \institution{Samsung Research}
  \country{Republic of Korea}
}
\email{{jieuns.shin, minsoo3.kim, kl.kim, james.geraci}@samsung.com}

%==========================================================
% Abstract
%==========================================================
\begin{abstract}
Proper data selection is essential for achieving high performance in
instruction tuning (IT) as the subset selected from a candidate pool
can substantially affect the capabilities acquired by the tuned
model. Recent Reinforcement Learning (RL)-based selectors formulate
data selection as a sequential decision problem and combine
training-time signals---loss-based difficulty $\mathcal{L}_i$ and
entropy-based uncertainty $\mathcal{H}_i$---into a composite reward
$\mathcal{R}_i = w\,\mathcal{L}_i + (1-w)\,\mathcal{H}_i$
with an adaptive variance-ratio weight~$w$.

We build on this RL\,+\,Composite Reward (CR) framework with
\textbf{Trajectory-Anchored Dynamic Selection (TADS)}. TADS extracts
a per-layer capability direction $v_l^{(t)}$ directly from the
model's own hidden-state trajectory, requiring no in-domain seed
set, labels, auxiliary model, or validation set. The resulting
trajectory anchor is incorporated multiplicatively into the
selection score and is controlled by a single hyperparameter
$\lambda$; setting $\lambda=0$ recovers the RL+CR baseline, yielding
a clean ablation of the anchor term.

We prove that, under a positive eigengap, the anchor sequence is
stable:
$\|v_l^{(t+1)} - v_l^{(t)}\| \le
(2C_\Sigma/\gamma_{\min})\,\eta_t$.
To our knowledge, this provides one of the first explicit stability
guarantees for a dynamic IT data-selection criterion.

Across model scales spanning $0.5$B to $14$B parameters, multiple
instruction-tuning datasets, and a diverse suite of evaluation
benchmarks, TADS consistently outperforms strong baselines while
adding less than $0.5\%$ wall-clock overhead over the RL+CR
baseline.
\end{abstract}

\keywords{Instruction tuning,
          Data selection,
          Reinforcement learning for data selection,
          Hidden-state representation,
          Large language model fine-tuning}

\maketitle

\end{document}

